\title{Stegte risnudler med sesam kylling}
\PrepTime{35}
\CookingTime{35}
\CookingTempe{}
\TypeCooking{Familievenlige måltider\index[nemlig]{Familievenlig!Stegte risnudler med sesam kylling}}
\NbPerson{4}
\Image{0 100 800 900}{images/nemlig/stegte-risnudler-m-sesam-kylling.png}

\begin{ingredient}
    \begin{main}
        \item 2 fed Hvidløg\index[ingredient]{Hvidløg}
        \item 300g Gulerødder\index[ingredient]{Gulerødder}
        \item 40g Ingefær\index[ingredient]{Ingefær}
        \item 1 Broccoli\index[ingredient]{Broccoli}
        \item 400g Hakket Kylling\index[ingredient]{Hakket Kylling}
        \item Salt\index[ingredient]{Salt}
        \item Peber\index[ingredient]{Peber}
        \item 30g Sesamfrø\index[ingredient]{Sesamfrø}
        \item Madolie\index[ingredient]{Olie}
        \item Smør\index[ingredient]{Smør}
        \item 200g Risnudler\index[ingredient]{Risnudler}
        \item 4 spsk. Soyasauce\index[ingredient]{Soyasauce}
        \item 4 spsk. Østerssauce\index[ingredient]{Østerssauce}
        \item 2 spsk. Sukker\index[ingredient]{Sukker}
        \item Vand\index[ingredient]{Vand}
    \end{main}
\end{ingredient}
\begin{recipe}
    \step{Pil og hak hvidløg fint. Gem halvdelen af hvidløget til de stegte nudler.}
    \step{Skyl og skræl gulerødderne. Riv halvdelen på den fine side af et rivejern og skær resten i tynde stave. Gem stavene til de stegte nudler.}
    \step{Skyl og evt. skræl ingefær og riv på den fine side af et rivejern. Stil til side til de stegte nudler.}
    \step{Skyl broccoli, og del i buketter. Snit buketterne i skiver (se tip!). Sæt til side til de stegte nudler.}
    \step{Kom hvidløg og revet gulerod i en skål sammen med hakket kylling, peber og rigeligt salt. Rør farsen grundigt, til den er godt sej. Stil til side.}
    \step{Varm en stor stegepande op og tørrist sesamfrø i et par minutter til de dufter, pas på de ikke brænder. Stil til side i en lille skål.}
    \step{Form farsen til små kyllingeboller med en våd teske og din håndflade.}
    \step{Varm en blanding af olie og smør op stegepanden og steg kyllingebollerne i ca. 4--5 min. ved middelvarme, til de er gyldne og gennemstegte.}
    \step{Drys med de ristede sesamfrø, vend godt rundt og sæt kyllingebollerne til side i en skål. Lad stegepanden stå, den skal du bruge om lidt.}
    \step{Kog nudler i letsaltet vand efter anvisningerne på pakken, til de næsten er færdigkogte men stadig har bid.}
    \step{Hæld nudlerne i en si og skyl dem under koldt vand. Lad dem dryppe af.}
    \step{Varm stegepanden op igen med lidt olie og steg hvidløg, ingefær, gulerodstave og broccoli i ca. 3--5 min.}
    \step{Tilsæt soyasauce, østerssauce og sukker og vend det hele godt rundt.}
    \step{Tilsæt nudlerne og smag til med ekstra soyasauce eller salt og peber. Bliver nudlerne for tørre, kan du med fordel tilsætte lidt vand.}
    \step{Servér de stegte risnudler i østersauce med sprøde grøntsager og top med sesam kyllingebollerne.}
\end{recipe}
\begin{notes}
    \textbf{Tip!} Du kan også sagtens spise broccoli stokken. Skær eller skræl det yderste af, det
kan være lidt trævlet. Snit så det inderste af stokken i skiver, og steg med i retten.
\end{notes}

\setcounter{step}{0}
\setcounter{substep}{0}

\clearpage